\documentclass[11pt]{article}
\usepackage{mathblatt}
\def\mitzone{1}
\begin{document}
\blattkopf{Flächen}{Gesamt}
\weit
\verzeichniszeile{\verz{zone}{Kennst du schon (Nr.~1–9)} \verztrenn \verz{e1}{1 Rechteck, Quadrat, Umfang (Nr.~10–14)} \verztrenn \verz{e2}{2 Parallelogramm (Nr.~15–17)} \verztrenn \verz{e3}{3 Dreieck (Nr.~18–21)} \verztrenn \verz{e4}{4 Trapez, Drachenviereck, Raute (Nr.~22–25)} \verztrenn \verz{e5}{5 Zusammengesetzte Figuren (Nr.~26–29)} \verztrenn \verz{abhaken}{Das kann ich}}
% Zone „Das kennst du schon“ – aus bank/flaechen/zone.jsonl (zusammenbau v0.1)
\einheitenkopf[zone][Kennst du schon]{Das kennst du schon}
\setcounter{aufgabe}{0}

%% TODO Titel: Ich-kann-Satz fehlt in der Bank (Platzhalter: Kettenname) – Nr. 1
%% TODO Anweisung: ein Satz über der ersten Teilaufgabe fehlt in der Bank – Nr. 1
\begin{aufgabe}{Längeneinheiten umrechnen (mm, cm, m), gemischte Angaben angleichen}
\begin{teile}
\teil $3$ m – wie viele cm? \leerfeld[cm]
\teil $1{,}35$ m – wie viele cm? \leerfeld[cm]
\end{teile}
\end{aufgabe}

%% TODO Titel: Ich-kann-Satz fehlt in der Bank (Platzhalter: Kettenname) – Nr. 2
%% TODO Anweisung: ein Satz über der ersten Teilaufgabe fehlt in der Bank – Nr. 2
\begin{aufgabe}{Flächeneinheiten (cm², m²) und Umrechnungszahl 100}
\begin{teile}
\teil $4$ cm² – wie viele mm²? \leerfeld mm²
\teil $250$ cm² – wie viele dm²? \leerfeld dm²
\end{teile}
\end{aufgabe}

%% TODO Titel: Ich-kann-Satz fehlt in der Bank (Platzhalter: Kettenname) – Nr. 3
%% TODO Anweisung: ein Satz über der ersten Teilaufgabe fehlt in der Bank – Nr. 3
\begin{aufgabe}{Multiplizieren und Dividieren mit Dezimalzahlen (Kommazahl mal Kommazahl, zweistellig geteilt durch einstellig)}
\begin{teile}
\teil $0{,}5 \cdot 8$ – wie viel? \leerfeld
\teil $2{,}5 \cdot 1{,}4$ – wie viel? \leerfeld
\end{teile}
\end{aufgabe}

%% TODO Titel: Ich-kann-Satz fehlt in der Bank (Platzhalter: Kettenname) – Nr. 4
%% TODO Anweisung: ein Satz über der ersten Teilaufgabe fehlt in der Bank – Nr. 4
\begin{aufgabe}{Formel nach einer Größe umstellen (A = a · b → b = A : a)}
\begin{teile}
\teil $P = q \cdot r$ – nach $r$ umstellen?
\teil $G = m \cdot p$ mit $G = 7{,}2$ und $m = 3$ – $p$? p = \leerfeld
\end{teile}
\end{aufgabe}

%% TODO Titel: Ich-kann-Satz fehlt in der Bank (Platzhalter: Kettenname) – Nr. 5
%% TODO Anweisung: ein Satz über der ersten Teilaufgabe fehlt in der Bank – Nr. 5
\begin{aufgabe}{Wurzel ziehen bei Quadratzahlen (√49)}
\begin{teile}
\teil $\sqrt{25}$ – wie viel? \leerfeld
\teil $\sqrt{144}$ – wie viel? \leerfeld
\end{teile}
\end{aufgabe}

%% TODO Titel: Ich-kann-Satz fehlt in der Bank (Platzhalter: Kettenname) – Nr. 6
%% TODO Anweisung: ein Satz über der ersten Teilaufgabe fehlt in der Bank – Nr. 6
\begin{aufgabe}{Rechten Winkel erkennen und einzeichnen (Geodreieck)}
\begin{teile}
\teil Ist das ein rechter Winkel? \janein

\winkel{90}{\alpha}
\teil An welcher Ecke ist ein rechter Winkel? Prüfe mit dem Geodreieck und markiere ihn.

\viereck{(0,0)}{(5,0)}{(5,3)}{(1,2.5)}
\end{teile}
\end{aufgabe}

%% TODO Titel: Ich-kann-Satz fehlt in der Bank (Platzhalter: Kettenname) – Nr. 7
%% TODO Anweisung: ein Satz über der ersten Teilaufgabe fehlt in der Bank – Nr. 7
\begin{aufgabe}{Kreisfläche (π · r²)}
\begin{teile}
\teil Kreis mit Radius $r = 1$ m – Fläche? Runde auf zwei Stellen. \leerfeld[m²]
\teil Kreis mit Radius $r = 1{,}5$ m – Fläche? Runde auf zwei Stellen. \leerfeld[m²]
\end{teile}
\end{aufgabe}

%% TODO Titel: Ich-kann-Satz fehlt in der Bank (Platzhalter: Kettenname) – Nr. 8
%% TODO Anweisung: ein Satz über der ersten Teilaufgabe fehlt in der Bank – Nr. 8
\begin{aufgabe}{Flächeneinheiten (cm², m²) und Umrechnungszahl 100 – Fehler finden}
\begin{teile}
\teil Mia rechnet: $8$ cm² $= 80$ mm². Wo ist der Fehler?
\end{teile}
\end{aufgabe}

%% TODO Titel: Ich-kann-Satz fehlt in der Bank (Platzhalter: Kettenname) – Nr. 9
%% TODO Anweisung: ein Satz über der ersten Teilaufgabe fehlt in der Bank – Nr. 9
\begin{aufgabe}{Flächeneinheiten (cm², m²) und Umrechnungszahl 100 – selbst rechnen}
\begin{teile}
\teil $9$ dm² – wie viele cm²? \leerfeld[cm²]
\end{teile}
\end{aufgabe}

\clearpage
% Einheit 1 von 5 – Rechteck, Quadrat, Umfang (bank/flaechen/e1.jsonl, zusammenbau v0.1)
%% TODO Zweigzeile Teil 1: was hier gelernt wird (Satz in Schülersprache aus dem Einheitstitel) – Einheit 1
\einheitenkopf[e1]{Einheit 1 von 5 · Rechteck, Quadrat, Umfang}
\zweigzeile{ab Kl. 5 · P10 oft}
\setcounter{aufgabe}{9}

%% TODO Titel: Ich-kann-Satz fehlt in der Bank (Platzhalter: Kettenname) – Nr. 10
%% TODO Anweisung: ein Satz über der ersten Teilaufgabe fehlt in der Bank – Nr. 10
\begin{aufgabe}{Gegeben – gesucht}
\begin{teile}
\teil Ein Rechteck hat die Fläche $A = 32$ m² und die Breite $b = 4$ m. Was ist gesucht? Kreise es in $A = a \cdot b$ ein.
\end{teile}
\end{aufgabe}

%% TODO Titel: Ich-kann-Satz fehlt in der Bank (Platzhalter: Kettenname) – Nr. 11
%% TODO Anweisung: ein Satz über der ersten Teilaufgabe fehlt in der Bank – Nr. 11
\begin{aufgabe}{Rechteck}
%% TODO \rechenplatz{n}: Schrittzahl des Grundfalls fehlt in der Bank (nur bei fester Schreibform, 2.3 b)
\begin{teile}
\teil Ein Garten bekommt einen Zaun. Fläche oder Umfang? \\ \kreuz{Fläche} \\ \kreuz{Umfang}
\teil Rechteck, $a = 8$ cm, $b = 3$ cm – Fläche und Umfang? A = \leerfeld[cm²] , u = \leerfeld[cm]
\teil Rechteck, $a = 4{,}5$ cm, $b = 2$ cm – Fläche und Umfang? A = \leerfeld[cm²] , u = \leerfeld[cm]
\teil Quadrat, $a = 6$ cm – Fläche und Umfang? A = \leerfeld[cm²] , u = \leerfeld[cm]
\teil Ein Fünfeck hat die Seiten $3$ cm, $4$ cm, $5$ cm, $2$ cm und $6$ cm – Umfang? u = \leerfeld[cm]
\teil Rechteck, $A = 56$ cm², $a = 8$ cm – wie lang ist $b$? b = \leerfeld[cm]
\teil Rechteck, $u = 20$ cm, $a = 6$ cm – wie lang ist $b$? b = \leerfeld[cm]
\teil Quadrat, $A = 64$ cm² – Seitenlänge? a = \leerfeld[cm]
\teil Ein Rechteck ist doppelt so lang wie breit; die Breite ist $x$. Term für den Umfang? u = \leerfeld
\teil Ein rechteckiger Bilderrahmen hat einen Umfang von $34$ cm, eine Seite ist $11$ cm lang. Wie lang ist die andere Seite? \hfill (P10 2018 OS) \\ \leerfeld[cm]
\end{teile}
\end{aufgabe}

%% TODO Titel: Ich-kann-Satz fehlt in der Bank (Platzhalter: Kettenname) – Nr. 12
%% TODO Anweisung: ein Satz über der ersten Teilaufgabe fehlt in der Bank – Nr. 12
\begin{aufgabe}{aus Rechtecken zusammengesetzte Fläche (Summe)}
\begin{teile}
\teil Ein L-förmiger Raum besteht aus zwei Rechtecken: $6$ m × $4$ m und $3$ m × $2$ m. Fläche? \leerfeld[m²]
\end{teile}
\end{aufgabe}

%% TODO Titel: Ich-kann-Satz fehlt in der Bank (Platzhalter: Kettenname) – Nr. 13
%% TODO Anweisung: ein Satz über der ersten Teilaufgabe fehlt in der Bank – Nr. 13
\begin{aufgabe}{Einheit wechseln vor dem Rechnen (cm und m gemischt)}
\begin{teile}
\teil Rechteck, $2$ m lang, $45$ cm breit – Fläche in m²? \leerfeld[m²]
\end{teile}
\end{aufgabe}

%% TODO Titel: Ich-kann-Satz fehlt in der Bank (Platzhalter: Kettenname) – Nr. 14
%% TODO Anweisung: ein Satz über der ersten Teilaufgabe fehlt in der Bank – Nr. 14
\begin{aufgabe}{Rechteck – Fehler finden, Begründen, Anwendung, Darstellungswechsel}
\begin{teile}
\teil Ein Rechteck ist $6$ cm lang und $2{,}5$ cm breit. Tom rechnet den Umfang: \rechnung{u &= 6 \cdot 2{,}5 \\ u &= 15 \text{ cm}} Wo ist der Fehler?
\teil Warum rechnet man beim Rechteck Länge mal Breite? Erkläre mit Einheitsquadraten.
\teil Ein Zimmer ist $4{,}2$ m lang und $3{,}5$ m breit. Laminat gibt es in Paketen zu $2{,}2$ m². Wie viele Pakete muss man kaufen?
\teil Zeichne ein Rechteck, dessen Fläche $A = 3 \cdot x$ ist, und beschrifte die Seiten.

\rechenplatz[halb]{4}
\end{teile}
\end{aufgabe}

\clearpage
% Einheit 2 von 5 – Parallelogramm (bank/flaechen/e2.jsonl, zusammenbau v0.1)
%% TODO Zweigzeile Teil 1: was hier gelernt wird (Satz in Schülersprache aus dem Einheitstitel) – Einheit 2
\einheitenkopf[e2]{Einheit 2 von 5 · Parallelogramm}
\zweigzeile{ab Kl. 7, je nach Buch bis Kl. 8 · keine P10-Aufgabe}
\setcounter{aufgabe}{14}

%% TODO Titel: Ich-kann-Satz fehlt in der Bank (Platzhalter: Kettenname) – Nr. 15
%% TODO Anweisung: ein Satz über der ersten Teilaufgabe fehlt in der Bank – Nr. 15
\begin{aufgabe}{Parallelogramm}
%% TODO \rechenplatz{n}: Schrittzahl des Grundfalls fehlt in der Bank (nur bei fester Schreibform, 2.3 b)
\begin{teile}
\teil Zeichne die Höhe zur Grundseite $g$ ein und markiere den rechten Winkel.

\parallelogramm[seiten={$g$,,,}]{4}{2.5}{60}
\teil Parallelogramm, $g = 9$ cm, $h = 4$ cm – Fläche? A = \leerfeld[cm²]
\teil Parallelogramm: Grundseite $9$ cm, schräge Seite $6$ cm, Höhe $5$ cm – Fläche? A = \leerfeld[cm²]
\teil Parallelogramm, $g = 6{,}5$ cm, $h = 4$ cm – Fläche? A = \leerfeld[cm²]
\teil Parallelogramm mit den Seiten $6$ cm und $4$ cm. Die Höhe zur $6$-cm-Seite ist $2$ cm, die Höhe zur $4$-cm-Seite $3$ cm. Rechne die Fläche auf beide Arten. \leerfeld[cm²] und \leerfeld[cm²]
\teil Parallelogramm, $A = 63$ cm², $g = 9$ cm – Höhe? h = \leerfeld[cm]
\teil Ein Parallelogramm hat die Grundseite $g = 6{,}4$ cm und die Höhe $h = 3{,}5$ cm. Zeichne in der Skizze die Höhe zu $g$ ein, berechne die Fläche und begründe die Formel: Wie wird aus dem Parallelogramm ein Rechteck?

\parallelogramm[seiten={$g$,,,}]{4.5}{2.5}{60}
\end{teile}
\end{aufgabe}

%% TODO Titel: Ich-kann-Satz fehlt in der Bank (Platzhalter: Kettenname) – Nr. 16
%% TODO Anweisung: ein Satz über der ersten Teilaufgabe fehlt in der Bank – Nr. 16
\begin{aufgabe}{Umfang aus den Seiten}
\begin{teile}
\teil Parallelogramm mit den Seiten $9$ cm und $6$ cm – Umfang? u = \leerfeld[cm]
\end{teile}
\end{aufgabe}

%% TODO Titel: Ich-kann-Satz fehlt in der Bank (Platzhalter: Kettenname) – Nr. 17
%% TODO Anweisung: ein Satz über der ersten Teilaufgabe fehlt in der Bank – Nr. 17
\begin{aufgabe}{Parallelogramm – Fehler finden, Begründen, Anwendung}
\begin{teile}
\teil Parallelogramm: Grundseite $7$ cm, schräge Seite $5$ cm, Höhe $4$ cm. Nils rechnet: \rechnung{A &= 7 \cdot 5 \\ A &= 35 \text{ cm}^2} Wo ist der Fehler?
\teil Warum ist die Fläche eines Parallelogramms Grundseite mal Höhe? Begründe mit Abschneiden und Anlegen.
\teil Ein Parkstreifen hat die Form eines Parallelogramms: $45$ m lang (Grundseite) und $5{,}5$ m breit (Höhe). Er wird gepflastert, $1$ m² kostet $38$ €. Was kostet das Pflaster?
\end{teile}
\end{aufgabe}

\clearpage
% Einheit 3 von 5 – Dreieck (bank/flaechen/e3.jsonl, zusammenbau v0.1)
%% TODO Zweigzeile Teil 1: was hier gelernt wird (Satz in Schülersprache aus dem Einheitstitel) – Einheit 3
\einheitenkopf[e3]{Einheit 3 von 5 · Dreieck}
\zweigzeile{ab Kl. 7, je nach Buch bis Kl. 8 (am Gymnasium ab Kl. 5) · P10 oft}
\setcounter{aufgabe}{17}

%% TODO Titel: Ich-kann-Satz fehlt in der Bank (Platzhalter: Kettenname) – Nr. 18
%% TODO Anweisung: ein Satz über der ersten Teilaufgabe fehlt in der Bank – Nr. 18
\begin{aufgabe}{Dreieck}
%% TODO \rechenplatz{n}: Schrittzahl des Grundfalls fehlt in der Bank (nur bei fester Schreibform, 2.3 b)
\begin{teile}
\teil Zeichne die Höhe zur Grundseite $g$ ein und markiere den rechten Winkel.

\dreieck{(0,0)}{(5,0)}{(2,3)}{}{}{$g$}{}{}{}
\teil Dreieck, $g = 8$ cm, $h = 7$ cm – Fläche? A = \leerfeld[cm²]
\teil Rechtwinkliges Dreieck – Fläche?

\dreieckrw{5}{2.67}{$15$ cm}{$8$ cm}{$17$ cm}
A = \leerfeld[cm²]
\teil Stumpfwinkliges Dreieck, Grundseite $6$ cm; die Höhe dazu liegt außerhalb und ist $4$ cm lang – Fläche?

\dreieck{(0,0)}{(3,0)}{(4.5,2)}{}{}{$6$ cm}{}{}{}
A = \leerfeld[cm²]
\teil Dreieck, $g = 4{,}6$ cm, $h = 3$ cm – Fläche? A = \leerfeld[cm²]
\teil Dreieck, $A = 26$ cm², $h = 4$ cm – Grundseite? g = \leerfeld[cm]
\teil Ein gleichseitiges Dreieck hat die Seitenlänge $s$. Term für den Umfang? u = \leerfeld
\teil Ein dreieckiges Grundstück ABC hat bei C einen rechten Winkel; $BC = 16$ m, $AC = 45$ m. Ein Teilstück QBC mit Q auf AC soll $120$ m² groß sein. Wie lang ist QC? \hfill (P10 2020 OS) \\ QC = \leerfeld[m]
\end{teile}
\end{aufgabe}

%% TODO Titel: Ich-kann-Satz fehlt in der Bank (Platzhalter: Kettenname) – Nr. 19
%% TODO Anweisung: ein Satz über der ersten Teilaufgabe fehlt in der Bank – Nr. 19
\begin{aufgabe}{Höhe zur passenden Grundseite wählen (drei Höhen)}
\begin{teile}
\teil Dreieck ABC: $AB = 8$ cm, $BC = 7{,}5$ cm. Die Höhe auf AB ist $6$ cm, die Höhe auf BC $6{,}4$ cm. Fläche – rechne mit BC und der passenden Höhe. A = \leerfeld[cm²]
\end{teile}
\end{aufgabe}

%% TODO Titel: Ich-kann-Satz fehlt in der Bank (Platzhalter: Kettenname) – Nr. 20
%% TODO Anweisung: ein Satz über der ersten Teilaufgabe fehlt in der Bank – Nr. 20
\begin{aufgabe}{Umfang}
\begin{teile}
\teil Dreieck mit den Seiten $5$ cm, $7$ cm und $9$ cm – Umfang? u = \leerfeld[cm]
\end{teile}
\end{aufgabe}

%% TODO Titel: Ich-kann-Satz fehlt in der Bank (Platzhalter: Kettenname) – Nr. 21
%% TODO Anweisung: ein Satz über der ersten Teilaufgabe fehlt in der Bank – Nr. 21
\begin{aufgabe}{Dreieck – Fehler finden, Begründen, Anwendung, Darstellungswechsel}
\begin{teile}
\teil Dreieck mit $g = 9$ cm und $h = 6$ cm. Jana rechnet: \rechnung{A &= 9 \cdot 6 \\ A &= 54 \text{ cm}^2} Wo ist der Fehler?
\teil Warum rechnet man beim Dreieck Grundseite mal Höhe und nimmt davon die Hälfte? Begründe mit einem Parallelogramm.
\teil Ein dreieckiger Hausgiebel ist unten $8{,}4$ m breit und $3{,}5$ m hoch. Er wird gestrichen; eine Dose Farbe reicht für $6$ m². Wie viele Dosen braucht man?
\teil Zeichne ein Dreieck, dessen Umfang $u = 3 \cdot k$ ist, und beschrifte die Seiten.

\rechenplatz[halb]{4}
\end{teile}
\end{aufgabe}

\clearpage
% Einheit 4 von 5 – Trapez, Drachenviereck, Raute (bank/flaechen/e4.jsonl, zusammenbau v0.1)
%% TODO Zweigzeile Teil 1: was hier gelernt wird (Satz in Schülersprache aus dem Einheitstitel) – Einheit 4
\einheitenkopf[e4]{Einheit 4 von 5 · Trapez, Drachenviereck, Raute}
\zweigzeile{ab Kl. 7, je nach Buch bis Kl. 8 · P10}
\setcounter{aufgabe}{21}

%% TODO Titel: Ich-kann-Satz fehlt in der Bank (Platzhalter: Kettenname) – Nr. 22
%% TODO Anweisung: ein Satz über der ersten Teilaufgabe fehlt in der Bank – Nr. 22
\begin{aufgabe}{Trapez}
%% TODO \rechenplatz{n}: Schrittzahl des Grundfalls fehlt in der Bank (nur bei fester Schreibform, 2.3 b)
\begin{teile}
\teil Welche Formel passt zur Figur? ($a$, $c$ parallele Seiten, $h$ Höhe, $e$, $f$ Diagonalen) \\ \kreuz{$A = a \cdot h$} \\ \kreuz{$A = \frac{a + c}{2} \cdot h$} \\ \kreuz{$A = \frac{e \cdot f}{2}$}

\trapez[hoehe=h]{5}{3}{2.5}
\teil Trapez, $a = 9$ cm, $c = 5$ cm, $h = 4$ cm – Fläche? A = \leerfeld[cm²]
\teil Trapez, $a = 7{,}5$ cm, $c = 4{,}5$ cm, $h = 3$ cm – Fläche? A = \leerfeld[cm²]
\teil Trapez, $A = 42$ cm², $a = 9$ cm, $c = 5$ cm – Höhe? h = \leerfeld[cm]
\teil Drachenviereck mit den Diagonalen $e = 7$ cm und $f = 6$ cm – Fläche? A = \leerfeld[cm²]
\teil Gleichschenkliges Trapez: $a = 11$ cm, $c = 5$ cm, Schenkel je $5$ cm, Höhe $4$ cm – Umfang? u = \leerfeld[cm]
\teil Ein Regalbrett hat die Form eines gleichschenkligen Trapezes mit den parallelen Seiten $90$ cm und $70$ cm; seine Fläche ist $2\,960$ cm². Passen Blumentöpfe mit $35$ cm Durchmesser darauf? Weise es über die Tiefe (Höhe) nach. \hfill (P10 2014 OS)
\end{teile}
\end{aufgabe}

%% TODO Titel: Ich-kann-Satz fehlt in der Bank (Platzhalter: Kettenname) – Nr. 23
%% TODO Anweisung: ein Satz über der ersten Teilaufgabe fehlt in der Bank – Nr. 23
\begin{aufgabe}{Umfang mit Schenkel aus Pythagoras oder Sinus (Vorrat)}
\begin{teile}
\teil Gleichschenkliges Trapez: $a = 14$ cm, $c = 8$ cm, $h = 4$ cm – Umfang? u = \leerfeld[cm]
\end{teile}
\end{aufgabe}

%% TODO Titel: Ich-kann-Satz fehlt in der Bank (Platzhalter: Kettenname) – Nr. 24
%% TODO Anweisung: ein Satz über der ersten Teilaufgabe fehlt in der Bank – Nr. 24
\begin{aufgabe}{Diagonale aus A}
\begin{teile}
\teil Drachenviereck, $A = 35$ cm², $e = 7$ cm – wie lang ist $f$? f = \leerfeld[cm]
\end{teile}
\end{aufgabe}

%% TODO Titel: Ich-kann-Satz fehlt in der Bank (Platzhalter: Kettenname) – Nr. 25
%% TODO Anweisung: ein Satz über der ersten Teilaufgabe fehlt in der Bank – Nr. 25
\begin{aufgabe}{Trapez – Fehler finden, Begründen, Anwendung}
\begin{teile}
\teil Trapez mit $a = 15$ cm, $c = 9$ cm, $h = 7$ cm. Leo rechnet: \rechnung{A &= 15 \cdot 7 \\ A &= 105 \text{ cm}^2} Wo ist der Fehler?
\teil Warum rechnet man beim Trapez „Mittelwert der parallelen Seiten mal Höhe“? Begründe mit zwei gleichen Trapezen.
\teil Ein Deich hat im Querschnitt die Form eines Trapezes: unten $32$ m, oben $6$ m breit, $7$ m hoch. Wie viel m³ Erde braucht man für $100$ m Deich?
\end{teile}
\end{aufgabe}

\clearpage
% Einheit 5 von 5 – Zusammengesetzte Figuren (bank/flaechen/e5.jsonl, zusammenbau v0.1)
%% TODO Zweigzeile Teil 1: was hier gelernt wird (Satz in Schülersprache aus dem Einheitstitel) – Einheit 5
\einheitenkopf[e5]{Einheit 5 von 5 · Zusammengesetzte Figuren}
\zweigzeile{ab Kl. 5, je nach Buch bis Kl. 6 · P10 oft}
\setcounter{aufgabe}{25}

%% TODO Titel: Ich-kann-Satz fehlt in der Bank (Platzhalter: Kettenname) – Nr. 26
%% TODO Anweisung: ein Satz über der ersten Teilaufgabe fehlt in der Bank – Nr. 26
\begin{aufgabe}{Zusammengesetzte Figuren}
%% TODO \rechenplatz{n}: Schrittzahl des Grundfalls fehlt in der Bank (nur bei fester Schreibform, 2.3 b)
\begin{teile}
\teil In welche Teilflächen zerlegst du die Figur? Benenne sie.

\viereck{(0,0)}{(6,0)}{(6,3)}{(2,3)}
\teil Ein L-förmiges Beet besteht aus zwei Rechtecken: $8$ m × $3$ m und $3$ m × $4$ m. Fläche? A = \leerfeld[m²]
\teil Eine Giebelwand: unten ein Rechteck, $8$ m breit und $3$ m hoch, darauf ein Dreieck mit derselben Grundseite und $2{,}5$ m Höhe. Fläche der Wand? A = \leerfeld[m²]
\teil Einem Rechteck $9$ m × $7$ m fehlt an einer Ecke ein Rechteck $4$ m × $3$ m. Fläche? A = \leerfeld[m²]
\teil Aus einer Holzplatte $70$ cm × $45$ cm wird ein Quadrat mit $15$ cm Seitenlänge ausgesägt. Restfläche? A = \leerfeld[cm²]
\teil Eine quadratische Tischplatte mit $90$ cm Seitenlänge hat in der Mitte ein rundes Loch mit $r = 5$ cm für den Sonnenschirm. Restfläche? Runde auf ganze cm². A = \leerfeld[cm²]
\teil Einem Rechteck $11$ m × $6$ m fehlt an einer Ecke ein Rechteck $5$ m × $2$ m. Umfang der Restfigur? u = \leerfeld[m]
\teil Eine Eisbahn besteht aus einem Rechteck $24$ m × $14$ m und an beiden kurzen Seiten je einem Halbkreis mit $14$ m Durchmesser. Berechne die Fläche der Eisbahn. Runde auf ganze m². \hfill (P10 2017 OS) \\ A = \leerfeld[m²]
\end{teile}
\end{aufgabe}

%% TODO Titel: Ich-kann-Satz fehlt in der Bank (Platzhalter: Kettenname) – Nr. 27
%% TODO Anweisung: ein Satz über der ersten Teilaufgabe fehlt in der Bank – Nr. 27
\begin{aufgabe}{Sachaufgabe mit Entscheidung (reicht die Platte?)}
\begin{teile}
\teil Für eine Tischplatte $1{,}4$ m × $0{,}75$ m hat Paul eine Holzplatte mit $1{,}2$ m² Fläche, die $0{,}8$ m breit ist. Reicht sie?
\end{teile}
\end{aufgabe}

%% TODO Titel: Ich-kann-Satz fehlt in der Bank (Platzhalter: Kettenname) – Nr. 28
%% TODO Anweisung: ein Satz über der ersten Teilaufgabe fehlt in der Bank – Nr. 28
\begin{aufgabe}{Verschnitt in Prozent (Vorrat, Niveau III)}
\begin{teile}
\teil Aus einem Brett $2$ m × $0{,}25$ m werden $6$ Stücke je $0{,}3$ m × $0{,}25$ m gesägt. Wie viel Prozent des Bretts sind Verschnitt?
\end{teile}
\end{aufgabe}

%% TODO Titel: Ich-kann-Satz fehlt in der Bank (Platzhalter: Kettenname) – Nr. 29
%% TODO Anweisung: ein Satz über der ersten Teilaufgabe fehlt in der Bank – Nr. 29
\begin{aufgabe}{Zusammengesetzte Figuren – Fehler finden, Begründen, Anwendung}
\begin{teile}
\teil Ein Rasen ist $15$ m × $8$ m groß, darin liegt ein runder Teich mit $r = 2$ m. Ina rechnet die Rasenfläche: \rechnung{A &= 15 \cdot 8 + \pi \cdot 2^2 \\ A &\approx 132{,}6 \text{ m}^2} Wo ist der Fehler?
\teil Eine L-Figur lässt sich auf zwei Arten in zwei Rechtecke zerlegen. Kommt beide Male dieselbe Fläche heraus? Begründe.
\teil Ein L-förmiges Zimmer besteht aus den Rechtecken $5$ m × $4$ m und $2{,}5$ m × $2$ m. Es bekommt Parkett für $32$ € je m²; für Verschnitt kauft man $5\,\%$ mehr. Was kostet das Parkett?
\end{teile}
\end{aufgabe}

% Abhakseite „Das kann ich“ – Zone nur im Gesamt (\mitzone)
%% TODO Ich-kann-Sätze: Platzhalter wie in den Aufgabendateien
\begin{abhakseite}
\ifdefined\mitzone
\abhakgruppe{Kennst du schon}
\abhak{1}{Längeneinheiten umrechnen (mm, cm, m), gemischte Angaben angleichen}
\abhak{2}{Flächeneinheiten (cm², m²) und Umrechnungszahl 100}
\abhak{3}{Multiplizieren und Dividieren mit Dezimalzahlen (Kommazahl mal Kommazahl, zweistellig geteilt durch einstellig)}
\abhak{4}{Formel nach einer Größe umstellen (A = a · b → b = A : a)}
\abhak{5}{Wurzel ziehen bei Quadratzahlen (√49)}
\abhak{6}{Rechten Winkel erkennen und einzeichnen (Geodreieck)}
\abhak{7}{Kreisfläche (π · r²)}
\abhak{8}{Flächeneinheiten (cm², m²) und Umrechnungszahl 100 – Fehler finden}
\abhak{9}{Flächeneinheiten (cm², m²) und Umrechnungszahl 100 – selbst rechnen}
\fi
\abhakgruppe{Einheit 1 · Rechteck, Quadrat, Umfang}
\abhak{10}{Gegeben – gesucht}
\abhak{11}{Rechteck}
\abhak{12}{aus Rechtecken zusammengesetzte Fläche (Summe)}
\abhak{13}{Einheit wechseln vor dem Rechnen (cm und m gemischt)}
\abhak{14}{Rechteck – Fehler finden, Begründen, Anwendung, Darstellungswechsel}
\abhakgruppe{Einheit 2 · Parallelogramm}
\abhak{15}{Parallelogramm}
\abhak{16}{Umfang aus den Seiten}
\abhak{17}{Parallelogramm – Fehler finden, Begründen, Anwendung}
\abhakgruppe{Einheit 3 · Dreieck}
\abhak{18}{Dreieck}
\abhak{19}{Höhe zur passenden Grundseite wählen (drei Höhen)}
\abhak{20}{Umfang}
\abhak{21}{Dreieck – Fehler finden, Begründen, Anwendung, Darstellungswechsel}
\abhakgruppe{Einheit 4 · Trapez, Drachenviereck, Raute}
\abhak{22}{Trapez}
\abhak{23}{Umfang mit Schenkel aus Pythagoras oder Sinus (Vorrat)}
\abhak{24}{Diagonale aus A}
\abhak{25}{Trapez – Fehler finden, Begründen, Anwendung}
\abhakgruppe{Einheit 5 · Zusammengesetzte Figuren}
\abhak{26}{Zusammengesetzte Figuren}
\abhak{27}{Sachaufgabe mit Entscheidung (reicht die Platte?)}
\abhak{28}{Verschnitt in Prozent (Vorrat, Niveau III)}
\abhak{29}{Zusammengesetzte Figuren – Fehler finden, Begründen, Anwendung}
\end{abhakseite}

\end{document}
